\documentclass[11pt]{article}
\usepackage[margin=1in]{geometry}

% Core packages
\usepackage{amsmath,amssymb}
\usepackage{booktabs}
\usepackage{graphicx}
\usepackage{tabularx}
\usepackage{xcolor}
\usepackage[numbers,sort&compress]{natbib}
\usepackage[hidelinks]{hyperref}
\usepackage{subcaption}
\usepackage{float}
\usepackage{placeins}
\usepackage[strings]{underscore}
\graphicspath{{figures/}}

% Paragraphs
\setlength{\parindent}{0pt}
\setlength{\parskip}{0.65\baselineskip}

\newcommand{\bigO}{\mathcal{O}}

\title{Bigger Bird: Context-Aware Sparse Attention for Efficient Long-Context Scaling}
\author{University of Toronto Machine Intelligence Student Team (UTMIST)}
\date{January 2026 -- Present}

\begin{document}
\maketitle

\begin{abstract}

% The quadratic complexity of dense attention makes it challenging to efficiently scale Transformers to long contexts. Here we introduce Bigger Bird, using a context-aware sparse attention mechanism. This model combines attention sinks, a sliding local window, and content-based routing of distant key--value positions, where routing scores are derived from mean causal attention probabilities over a late-prompt query set. Our Triton kernel implements 64 sink tokens, a
% 16{,}384-token local window, a per-head routed budget of up to 512 positions, and exact dense attention over the final 256 prompt queries. We evaluate the
% method under a paired dense baseline on BF16 R1-Distill-Llama-8B across
% 1{,}178 long-context examples from LongCodeQA, Code.Debug, and LongBench~v2.
% Pooled accuracy is 34.6\% for Bigger Bird versus 35.3\% for dense attention,
% while mean prefill throughput improves by $1.09$--$1.27\times$ on the longer
% subsets. A complementary no-KV-cache LongCodeQA pilot at 542K--878K tokens
% yields a $10.23\times$ mean prefill speedup on 18 prompts. Together, these results suggest that attention routing can substantially reduce long-context prefill cost while keeping pooled accuracy close to dense attention. This points to a possible path toward more efficient long-context inference and motivates exploring whether sparse attention can also reduce the cost of training at longer contexts. However, our experiments evaluate inference only, and the small 1M-token pilot does not yet establish accuracy or efficiency at that scale. We include early IMDb and kernel experiments to show scale at different context windows, how early scale behaviour and the iteration process led up to the current method with design rationale. We focus on inference-time evaluation, and sparse-attention training is left as future work. 






The quadratic complexity of dense attention makes it challenging to efficiently scale Transformers to long contexts. Here we introduce Bigger Bird, using context-sensitive sparse attention. This model uses attention sinks, a sliding window, and routing of distant key-value pairs based on the average causal attention probabilities across a set of late prompt queries. In the Triton implementation, 64 attention sinks, 16,384 tokens within the sliding window, head-specific routed attention to up to 512 positions and exact dense attention on the final 256 prompt queries can be applied. We test the performance of this model against a dense baseline on BF16 R1-Distill-Llama-8B on 1,178 long context examples. Pooled accuracy was 34.6\% for Bigger Bird compared to 35.3\% using dense attention. Prefill speeds were improved by $1.09$-$1.27$x faster in the longest subsets of prompts. For further analysis, we performed a pilot experiment without using KV caches on LongCodeQA with prompt lengths between 542K--878K tokens. Using this data, we improve prefill speeds by a factor of $10.23$x averaged over 18 prompts. These findings indicate that sparse attention routing can significantly reduce prefill costs in long contexts, while maintaining pooled accuracy similar to dense attention. This suggests potential for more efficient long-context inference. In this work, we only consider inference performance, so it remains to be seen if we can also reduce training costs in long contexts using sparse attention. Furthermore, the pilot only covers up to 1M tokens, meaning that we cannot be certain of the accuracy or efficiency at this scale. Initial IMDb and attention-kernel experiments are shown to demonstrate scaling for different context window sizes, and document the iterative development and design rationale behind Bigger Bird. 


\end{abstract}

\newpage
\tableofcontents

\newpage
\section{Introduction}

% Transformer models rely on self-attention to connect each token with every other token in a sequence. Dense self-attention forms an $n\times n$ attention matrix, giving it $\bigO(n^2)$ comparison and memory cost in sequence length $n$ \citep{vaswani2017attention}. This quadratic dependence is a practical bottleneck for long-context inference and for any future training regime that must process long documents end to end.

% Bigger Bird investigates context-aware sparse attention as an alternative to paying this cost uniformly. Rather than committing to a fixed sparse graph, the method allocates computation to early attention sinks, a recent local window, and a content-selected set of distant key--value positions. The current system of record is \textbf{Bigger Bird v3}: a training-free Triton sparse-prefill kernel for decoder LLMs, evaluated under a paired dense baseline on R1-Distill-Llama-8B. Decode remains exact on the retained KV cache; optional KV eviction and related memory savers are studied as complementary inference-time mechanisms.

% The project developed through a broader architecture search. Experiments~0--15
% cover encoder-scale prototypes on BART / BigBird-RoBERTa \citep{zaheer2020bigbird}; Experiments~16--19 extend the sparse-attention search on R1; Experiment~20 is the real long-context benchmark suite; and Experiment~21 studies routing reuse and cache variants. After situating the work in the literature (Section~\ref{sec:related}), this draft presents the Bigger Bird v3 method, the experimental setup, the primary R1 results, and the developmental encoder evidence that explains how the project arrived at the current design.




The self-attention mechanism in Transformers provides pairwise interactions
between all tokens in a sequence. However, dense attention has
$\mathcal{O}(n^2)$ time and space complexity, where $n$ is the sequence
length \citep{vaswani2017attention}. This quadratic complexity limits
inference on long contexts and future end-to-end training on long documents.

To reduce this cost, Bigger Bird implements context-aware sparse
attention. The model computes attention weights over (i) early attention
targets, (ii) local attention windows, and (iii) a set of key--value pairs
determined by content.

This work presents a training-free sparse-prefill kernel for decoder LLMs
(Bigger Bird v3). We compare it against a dense baseline using
R1-Distill-Llama-8B. Full decoding uses the preserved KV cache. We also
investigate optional KV-cache eviction and other memory-efficient strategies
as complementary approaches for inference.

The experiments were part of a larger architecture search. Experiments
0--15 apply encoder-scale models to BART and BigBird-RoBERTa
\citep{zaheer2020bigbird}. Experiments 16--19 extend the sparse-attention
search to R1. Experiment 20 is a long-context benchmark suite, and
Experiment 21 investigates routing reuse and KV-cache variants. Following
the background (Section~\ref{sec:related}), we detail the methods used in Bigger Bird v3, our experimental setup, key results for R1 and provide evidence from encoder models that illustrates how we converged on our current design.

\section{Related Work}
\label{sec:related}

Bigger Bird~v3 combines three ideas that have appeared separately in prior
work: a fixed scaffold of sink and local tokens, long-range positions chosen
from the prompt's own attention, and a training-free kernel that applies this
selection during prefill. We review each line of work in turn, from fixed
patterns to the sparse-prefill methods closest to ours, and close with the
techniques used in our earlier encoder prototypes.

\paragraph{Fixed sparse patterns.}
Sparse Transformers \citep{child2019sparse} restrict each query to strided
and fixed patterns, reducing attention cost to $\bigO(n\sqrt{n})$. Longformer
\citep{beltagy2020longformer} and BigBird \citep{zaheer2020bigbird} combine a
sliding window with a small number of global tokens, and BigBird adds random
links. These patterns are chosen before the input is seen and are built into
the model during pretraining. Bigger Bird keeps their local-plus-global
structure, but chooses its long-range positions for each prompt and applies
them to a pretrained model without retraining.

\paragraph{Attention sinks and streaming.}
\citet{xiao2024streamingllm} observed that pretrained language models direct
a large share of attention to the first few tokens even when those tokens
carry little meaning, and that keeping these ``attention sinks'' alongside a
recent window lets a model run on long streams without fine-tuning. Our sink
and local components follow the same design. However, because StreamingLLM
discards everything outside the window, evidence in the middle of a long
prompt becomes unreachable; the routed positions in Bigger Bird are intended
to recover it.

\paragraph{KV-cache selection.}
A related line of work reduces decoding cost by keeping only part of the KV
cache. H2O \citep{zhang2023h2o} retains entries with high accumulated
attention, and SnapKV \citep{li2024snapkv} scores prompt positions for each
head using an observation window at the end of the prompt. Our router uses a
similar signal, the mean attention from the final 64 prompt queries, but
applies it to prefill instead: the selected positions determine which keys
each query attends to, while the full cache is kept for decoding.

\paragraph{Sparse prefill.}
The closest work makes prefill itself sparse. MInference
\citep{jiang2024minference} assigns each head a sparse pattern family offline
and estimates its indices for each prompt. DeepSeek Sparse Attention
\citep{deepseek2025v32} adds a lightweight indexer that selects the top-$k$
entries for each query; the indexer still scores the full preceding context,
and the model is adapted to the sparse pattern through continued training.
Bigger Bird differs from both by applying one pattern to every head, with no
offline pattern search and no training. More recent hierarchical routers, LSA
and ZHA, reduce routing cost further using key summaries; we have not
evaluated them and treat them as future comparisons.

\paragraph{Trainable sparse attention.}
Native Sparse Attention \citep{yuan2025nsa} combines compressed, selected,
and sliding-window branches, and MoBA \citep{lu2025moba} routes each query to
a few key blocks. Both learn their sparsity during pretraining or continued
training. They show that trained sparse attention can work at scale, but they
require changing the model's weights, so they are not directly comparable to
a training-free method. Our decoder screens included an untrained NSA-style
variant, which failed retrieval (Section~\ref{sec:results}).

\paragraph{Techniques from our encoder prototypes.}
The earlier prototypes in Section~\ref{sec:developmental} drew on a wider
range of ideas. Reformer \citep{kitaev2020reformer} and the Routing
Transformer \citep{roy2021routing} group similar queries and keys by hashing
or clustering; Experiments~1, 4, 6, and~10 instead score keys directly and
keep the top-$k$. PoWER-BERT \citep{goyal2020powerbert} and learned token
pruning \citep{kim2021tokenpruning} shorten the sequence as depth increases,
as do Token Dropping and Dynamic Context (Experiments~8 and~13). The encoder
Bigger Bird designs (Experiments~5 and~15) selected keys with maximal
marginal relevance \citep{carbonell1998mmr} and facility location
\citep{lin2011submodular}, two standard objectives for choosing a relevant
but non-redundant subset.

\section{Method}
\label{sec:methodology}
\label{sec:method}

This section states the current Bigger Bird v3 design first, then summarizes
the architecture search that produced it. Evaluation suites, protocols, and
metrics are in Section~\ref{sec:setup}.

\subsection{Bigger Bird v3}
\label{exp:bb-v3}
\label{sec:method-v3}

Bigger Bird v3 is the current method of record: a training-free sparse-prefill
attention for decoder LLMs. The design goal is to preserve useful long-range
links while avoiding a full causal attention matrix over the prompt.

Default configuration:
\texttt{\{"sink":64, "window":16384, "vertical":512, "exact\_tail":256\}}.
For each layer, every query attends to
\begin{itemize}
  \item the first 64 tokens (attention sinks);
  \item the previous 16{,}384 tokens (local window);
  \item up to 512 distant positions per head, selected from mean causal
    attention probability over the final 64 prompt queries (content-based
    routing);
  \item and, for the final 256 prompt queries, exact attention over their
    visible prefix (exact tail).
\end{itemize}
Decode is exact on the retained KV cache. Optional KV eviction and memory
savers (chunked MLP, CPU embeddings, INT8 KV) are complementary inference
mechanisms, not part of the core attention pattern. The implementation lives
in \texttt{modal\_rd/attention.py} with the Triton kernel in
\texttt{kernels/bigger\_bird\_flash.py}.


\subsection{Prior architecture search}
\label{sec:architectures}

Table~\ref{tab:taxonomy} summarizes the broader search that led to Bigger Bird
v3. Experiments~0--15 are encoder-scale prototypes; Experiments~16--21 are the
later R1 / long-context line. Brief mechanism notes follow; measured outcomes
appear in Sections~\ref{sec:results} and~\ref{sec:developmental}.

\begin{table}[ht]
  \centering
  \caption{Taxonomy of evaluated attention mechanisms (Experiments~0--21 and
  Bigger Bird v3).}
  \label{tab:taxonomy}
  \begin{tabularx}{\textwidth}{@{}l>{\raggedright\arraybackslash}Xl@{}}
    \toprule
    Family & Representative mechanism & Experiments \\
    \midrule
    Reference & Dense self-attention (BART; later HF SDPA on R1) & 0 \\
    Routed sparse & Low-rank top-$k$, block routing, GQA routing,
      coarse-to-fine & 1, 4, 6, 10, 17 \\
    Adaptive computation & Layer schedules, token dropping, dynamic context
      & 7, 8, 13, 14 \\
    Hybrid / multi-branch & Windows, globals, speculation, NSA, S2-HHST,
      free NSA, confidence-gated & 2, 3, 9, 11, 12, 16, 18 \\
    Encoder Bigger Bird & MMR / facility-location / teleports on BART or
      BigBird-RoBERTa & 5, 15 \\
    Decoder Bigger Bird & Flash / Triton sparse prefill; sinks + window +
      routed distant keys + exact tail & 19, v3 \\
    Long-context systems & Real-benchmark harness; routing reuse / KV
      eviction variants & 20, 21 \\
    \bottomrule
  \end{tabularx}
\end{table}

Experiments~11 and~12 began as unfinished prototypes. Experiment~20 is the
real-benchmark harness used in Section~\ref{sec:results}; Experiment~21 studies
routing reuse and KV-cache variants.


\subsubsection*{Experiment 0: Dense Baseline}
\label{exp:baseline}
Standard dense self-attention, every query attending to every key:
$\mathrm{Attention}(Q,K,V) = \mathrm{softmax}(QK^\top/\sqrt{d})V$. Early
encoder runs used BART \citep{lewis2020bart}; the current R1 baseline is
Hugging Face SDPA on R1-Distill-Llama-8B. This remains the reference for
quality and resource comparisons. \emph{Complexity:} $\bigO(n^2)$ compute and
memory in sequence length $n$.

\subsubsection*{Experiment 1: DeepSeek Top-k Routing}
\label{exp:deepseek}
A custom head-shared, low-rank router first scores all keys cheaply in a
reduced dimension, then selects a fixed number of top-scoring keys per head;
each query attends only to that shared key set under full-dimensional
attention. The defaults are $d_{\mathrm{low}}=16$ and top-$k=64$, with the
same keys reused across query positions in each head. The method falls back to
dense attention on short sequences. Although named after DeepSeek, it is an
adaptation rather than a direct implementation of Multi-head Latent Attention
\citep{deepseek2024v3}. \emph{Complexity:}
$\bigO(n\,d_{\mathrm{low}}+nKd)\approx\bigO(n)$ for fixed $K$.

\subsubsection*{Experiment 2: Lightning Hybrid}
\label{exp:lightning}
Combines an unfolded sliding-window local softmax, without materializing the
full $T\times T$ score matrix, with an optional linear-attention global path
using the feature map $\mathrm{elu}(x)+1$ \citep{qin2024lightning}. The global
path activates when the sequence exceeds four times the default 128-token
window, and its output is added to the local contribution.
\emph{Complexity:} $\bigO(nW)$ locally and $\bigO(nW+nd^2)$ with the global
path.

\subsubsection*{Experiment 3: Dynamic Global Tokens}
\label{exp:dynglobal}
A learned linear gate scores every token's ``globalness''; the top-scoring
tokens are selected once per batch as global tokens, and every query attends
to this learned global set plus a local sliding window in a single softmax.
\emph{Defaults:} $G=16$ global tokens, window size $W=64$.
\emph{Complexity:} $\bigO(n\,(G+W)\,d) = \bigO(n)$ for fixed $G,W$.

\subsubsection*{Experiment 4: PBS Attention}
\label{exp:pbs}
Queries and keys are mean-pooled into blocks, a head-shared block router
selects the top-scoring blocks, and tokens from those blocks are gathered with
sorted indices for coalesced GPU access before sparse softmax. The default is
two 64-token blocks, or 128 selected keys. \emph{Complexity:}
$\bigO(n_{\mathrm{blocks}}B+nMBd)\approx\bigO(n)$ for fixed block size $B$
and selected-block count $M$.

\subsubsection*{Experiment 5: Unified Bigger Bird (BART backbone)}
\label{exp:unified-bb}
Patches BART self-attention with one combined selection pattern: MMR-lite
diversity-aware local top-$k$ selection \citep{carbonell1998mmr}, a learned
gate that chooses global tokens, and biased random teleports that mix
high-gate candidates with uniformly sampled positions. Each query attends
only to the union of local selections, globals, and teleports. The learned
global gate deliberately differs from Experiment~15, which uses a submodular
facility-location objective.

\subsubsection*{Experiment 6: DeepSeek + PBS Hybrid}
\label{exp:deepseekpbs}
Combines content-aware routing with GPU-friendly block coalescing: low-rank
query/key projections (as in Experiment~1) first identify the top-scoring
blocks (as in Experiment~4), and top-$k$ tokens are then selected \emph{within}
those blocks; selected indices are sorted to guarantee contiguous memory
reads.

\subsubsection*{Experiment 7: Layer-Adaptive Sparsity}
\label{exp:layer-adaptive}
Uses a per-layer key budget that begins with broad attention and becomes more
sparse in later layers, reflecting the hypothesis that early layers require
more context for feature extraction while later layers can reason over fewer
selected tokens. The default budgets are 192 keys in early layers, 64 in
middle layers, and 32 in late layers.

\subsubsection*{Experiment 8: Token Dropping}
\label{exp:token-drop}
Runs three ordinary dense attention layers to extract local features, scores
each token's importance by the $\ell_2$ norm of its hidden state, and drops
the bottom 30\% of tokens before continuing with a shorter sequence for the
remaining layers. \emph{Complexity:} after the drop point, attention cost
scales with $\mathrm{keep\_ratio}^2$ while projection and FFN layers scale
linearly with $\mathrm{keep\_ratio}$ relative to the undropped sequence. This
follows the broader token-pruning literature \citep{goyal2020powerbert,kim2021tokenpruning}.

\subsubsection*{Experiment 9: Attention Speculation}
\label{exp:speculation}
Inspired by speculative decoding \citep{leviathan2023speculative}. Every layer runs a sparse ``fast path''
(local window plus first/middle/last anchor tokens); every fourth layer
additionally computes full dense attention as a verifier and adds a
KL-divergence distillation loss that trains the fast path to mimic full
attention. At inference, only the fast path executes.

The core complexity still remains  $\bigO(N^2)$ however, since the attention is pre-computed, it accelerates inference greatly. However, the algorithmic time complexity remains unchanged. Wrong speculation also has the same fallback 

\subsubsection*{Experiment 10: Grouped-Query Sparse Attention}
\label{exp:gqa}
Stacks two memory-saving mechanisms: (a) grouped-query attention compresses
12 key/value heads into 4 shared groups via mean pooling (a 3$\times$ memory
reduction), and (b) top-$k$ sparse routing is applied on top of the shared
groups, so softmax cost drops from $\bigO(N)$ to $\bigO(K)$ per query
\citep{ainslie2023gqa}.

\subsubsection*{Experiment 11: Native Sparse Attention (NSA)}
\label{exp:nsa}
This architecture combines three parallel attention branches per
query : coarse-grained compressed (pooled) attention, fine-grained
explicitly selected-block attention, and local sliding-window attention --
whose outputs are fused via learned per-query gates. This experiment is
marked incomplete (``TODO'') in the codebase at the time of this report; see
the note at the start of Section~\ref{sec:architectures} and the original NSA
design \citep{yuan2025nsa}. \\ 
\emph{Claims of the project:} Achieves sub-quadratic, hardware-aligned
attention complexity by operating on contiguous blocks rather than scattered
tokens, targeting significant reductions in memory bandwidth and FLOPs while
remaining end-to-end trainable and competitive with dense attention on long
sequences.\\ 
Classic dense attention
scales as $\bigO(N^2)$ and spends most of its compute attending to low-weight
distant tokens at full resolution. NSA separates global gist (compressed
branch), sharp long-range retrieval (selected-block branch), and local
context (sliding-window branch), preserving full-sequence information flow
at a fraction of the memory and compute cost.

\subsubsection*{Experiment 12: S2-HHST}
\label{exp:s2hhst}
Shard-based hierarchical sparse attention with
dedicated sink tokens: the sequence is partitioned into shards, attention is
computed hierarchically within each shard and across shard-level
representations, and a fixed set of sink tokens remains globally visible to
all queries. This experiment is likewise marked incomplete (``TODO'') in the
codebase; see the note at the start of Section~\ref{sec:architectures}.\\
\emph{Claims of the project:} Replaces the global $\bigO(N^2)$ attention
matrix with hierarchical shard-level scaling (local intra-shard cost plus
reduced cross-shard cost), lowering peak activation memory and enabling
longer context lengths than dense attention under the same hardware budget.\\ 
Classic attention forces
every token to attend directly to every other token, creating a quadratic
bottleneck. Hierarchical sharding avoids all-to-all token interactions while
still allowing information to propagate across the full sequence, and
dedicated sink tokens absorb excess softmax probability mass to keep sparse
attention distributions stable.

\subsubsection*{Experiment 13: Dynamic Context}
\label{exp:dynamic-context}
Extends Token Dropping to arbitrarily long inputs. After three dense early
layers, tokens are scored by hidden-state $\ell_2$ norm and up to 4096 are
retained regardless of the original length. Inputs longer than 8192 tokens
are processed independently by chunk in the early layers, avoiding a full
input-length attention matrix; late layers always attend over the fixed
budget. \emph{Complexity:}
$\bigO(\mathrm{num\_chunks}\times\mathrm{chunk\_size}^2) = \bigO(N \cdot \mathrm{chunk\_size})$
for early layers plus $\bigO(\mathrm{budget}^2)$ for late layers.

\subsubsection*{Experiment 14: Token Drop + DeepSeek Routing}
\label{exp:token-drop-deepseek}
 Combines Token Dropping with a DeepSeek-style
low-rank router for two levels of sparsity. Early layers run densely to
extract local features, 30\% of low-importance tokens (scored by hidden-state
$\ell_2$ norm) are dropped after layer 3, and subsequent layers apply
low-rank top-$k$ routing to the shortened sequence. Defaults are top-$k=64$
and low-rank dimension 16.\\
\emph{Claims of the project:} Compounds two reductions in late layers:
dropping tokens shrinks the active sequence length to $N' = 0.7N$ (reducing
both projection and FFN compute linearly), while top-$k$ routing reduces
late-layer value aggregation to $\bigO(N' \cdot k)$ using a cheap low-rank
scoring step ($\bigO(N'^2 d_r)$ with $d_r = 16 \ll d$).\\ Classic attention
allocates uniform compute to every token and every token pair across all
layers. By pruning redundant tokens early and restricting the surviving
tokens to their top-$k$ most relevant keys, this hybrid eliminates both
token-level redundancy and pairwise attention waste while preserving dense
feature extraction in the early layers.

\subsubsection*{Experiment 15: Proper Bigger Bird (BigBird-RoBERTa backbone)}
\label{exp:proper-bb}
The original UTMIST Bigger Bird design is built directly on
\texttt{bigbird-roberta-base} \citep{zaheer2020bigbird}. It extends the
middle-block attention path with MMR-lite local top-$k$ selection,
facility-location global selection \citep{lin2011submodular}, and biased
block-wise teleports. Defaults are 128-token fragments, $\max_k=64$, six
globals, and four teleports per head. No standalone task-level run is
available for this exact configuration; the related ablation below uses the
same selection ideas on a BART backbone.

\subsubsection*{BigBird-Family Ablation}
\label{exp:bigbird-family}
This separate study compares stock BigBird \citep{zaheer2020bigbird} with a
BART-based hybrid using blocked-MMR local selection, facility-location
globals, and optional teleports. It evaluates globals alone, local MMR alone,
a refined MMR rule, and combinations of MMR, globals, random links, and
teleports. Because its backbone differs from Experiment~15 and its global
selector differs from Experiment~5, the results isolate selection mechanisms
rather than directly compare the numbered Bigger Bird implementations.

\subsubsection*{Experiment 16: Free NSA}
Parameter-free / untrained NSA-style multi-branch routing (window +
compressed + selected). Used as a zero-shot R1 screen; learned compressors
and gates that are never trained fail retrieval in this setting.

\subsubsection*{Experiment 17: Coarse-to-Fine Routing}
Two-stage content routing: select top blocks, then top-$k$ tokens within
those blocks. Screens show that fine-token budget dominates block coverage at
long context.

\subsubsection*{Experiment 18: Confidence-Gated Multi-Resolution}
Local window plus a low-rank causal linear branch and a head-shared routed
global correction. Validation remains preliminary relative to Bigger Bird v3.

\subsubsection*{Experiment 19: Bigger Bird Flash}
An earlier Flash / Triton sparse-prefill prototype that preceded the final
v3 configuration. Retained in screening tables; superseded by Bigger Bird v3
for the main R1 claims.


\subsubsection*{Experiment 20: Real Long-Context Benchmark Harness}
Paired dense / Bigger Bird evaluation on LongCodeQA, Code.Debug, and
LongBench~v2. The full H100 matched suite (1{,}178 examples, 130K prompt cap)
is the primary real-benchmark evidence in Section~\ref{sec:r1-benchmark}.

\subsubsection*{Experiment 21: Kascade / Cache Variants}
Follow-up study of routing reuse across layers and KV-cache eviction /
compression at long context, including held-out fact retention probes near
1M tokens. Cross-layer route reuse is incompatible with per-layer eviction
when cache slots no longer refer to the same tokens.
\newpage
\section{Experimental Setup}
\label{sec:setup}

\subsection{Objectives}
\label{sec:objectives}

The evaluation addresses four objectives. The first two support the main
claims; the last two indicates how the final design was achieved.
\begin{enumerate}
  \item Test whether \textbf{context-aware sparse prefill preserves dense answer
    quality} on real long-context benchmarks (Table~\ref{tab:r1-benchmark}).
  \item Measure \textbf{prefill time, decode speed, and peak memory} under a matched
    dense baseline into the 32K--130K regime, with pilots beyond 500K tokens
    (Tables~\ref{tab:r1-benchmark}--\ref{tab:1m-evals}).
  \item Identify \textbf{which routing mechanisms retain near-dense accuracy} in
    synthetic retrieval screens
    (Tables~\ref{tab:screening}--\ref{tab:v3-ablate}).
  \item Characterize how \textbf{attention-kernel time, latency, and memory scale
    with sequence length} for early encoder-scale variants on IMDb reviews
    (Section~\ref{sec:developmental}).
\end{enumerate}

\subsection{Evaluation suites}
\label{sec:dataset}

\textbf{Primary (current).} The main quantitative claims use \texttt{BF16
R1-Distill-Llama-8B}~\citep{Guo_2025} under a paired dense baseline on three
real long-context benchmarks:
\begin{itemize}
  \item \textbf{LongCodeQA}~\citep{rando2025longcodebenchevaluatingcodingllms}
    repository comprehension (32K / 64K / 128K buckets, plus a 1M-token pilot);
  \item \textbf{InfiniteBench Code.Debug}~\citep{zhang-etal-2024-bench} bug
    localization;
  \item \textbf{LongBench~v2}~\citep{bai-etal-2025-longbench} multi-domain
    long-context reasoning.
\end{itemize}
Separately, synthetic \textbf{RULER} needle-in-a-haystack
tests~\citep{hsieh2024ruler} were used during method development to analyze
candidate designs, and they are not part of the main claims.

Table~\ref{tab:suites} summarizes the three real suites. All three are
four-way multiple choice, so chance accuracy is 25\%. Code.Debug and
LongBench~v2 are hard enough that even dense attention scores close to
chance (Section~\ref{sec:results}). Scoring and timing details are given in
Section~\ref{sec:protocol}.

\begin{table}[ht]
  \centering
  \small
  \caption{Long-context evaluation suites. All items are four-way multiple
    choice (chance $=25\%$). Prompts are capped at 130K tokens and
    middle-truncated beyond the cap. Mean tokens is the prompt length
    after truncation.}
  \label{tab:suites}
  \begin{tabular}{@{}lllrrr@{}}
    \toprule
    Suite & Task & Bucket & $n$ & Mean tokens & Truncated \\
    \midrule
    LongCodeQA~\citep{rando2025longcodebenchevaluatingcodingllms}
      & Repository QA & 32K  & 113 & $\sim$20K  & 0 \\
      &               & 64K  & 76  & $\sim$50K  & 0 \\
      &               & 128K & 92  & $\sim$88K  & 0 \\
    Code.Debug~\citep{zhang-etal-2024-bench}
      & Bug localization & -- & 394 & $\sim$108K & 136 \\
    LongBench~v2~\citep{bai-etal-2025-longbench}
      & Multi-domain reasoning & -- & 503 & $\sim$85K  & 186 \\
    \midrule
    Total & & & 1{,}178 & & 322 \\
    \bottomrule
  \end{tabular}
\end{table}

\textbf{Developmental (supporting).} Earlier work used \textbf{IMDb}
sentiment classification on BART / BigBird-RoBERTa at fixed lengths
\citep{maas2011imdb}, plus exploratory \textbf{Long Range Arena (LRA)}
ListOps and Text runs \citep{tay2021lra}. Most LRA runs did not rise above
chance or majority-class accuracy, so they are reported only in the appendix.
Those smaller-scale experiments document how the project searched
architectures, observed early scaling trends, and chose content-aware
routing. They explain how we arrived at Bigger Bird v3; they are not the
primary long-context claim.

\subsection{Experimental Protocol}
\label{sec:protocol}

\paragraph{Current R1 protocol.}
Primary runs use \texttt{BF16 R1-Distill-Llama-8B} on a single H100 80~GB with batch
size 1, comparing a paired dense baseline (HF SDPA under PyTorch's default
backend dispatch, which selects the FlashAttention kernel for causal BF16
inputs on the H100) against Bigger Bird v3. Both arms see identical prompts.
Prompts are capped at 130{,}000 tokens; longer prompts are middle-truncated,
because the task instructions and the question sit at the start and end of
each prompt. Prefill is a full-prompt forward with MLP chunking (8{,}192
tokens). Before timing, each arm runs two warmup forwards on synthetic
20{,}000- and 20{,}001-token prompts so that the Triton kernels are compiled
and specialized. Each item is then prefilled once per arm, and the reported
prefill times are means over items, not over repeated runs of the same item.
Decode timing uses a separate 16-step greedy CUDA-graph probe; it measures
speed only, not the quality of generated text or patch success.

Accuracy is read from the first answer-token distribution, renormalized over
the four options A--D. We score without chain-of-thought generation,
following the no-CoT setting of LongBench~v2: the prompt ends with an empty
think block and the cue ``The correct answer is ('', so the answer letter
depends only on the prefill. This isolates the effect of sparse prefill from
the variance of long reasoning traces and keeps the runs within the compute
budget. The matched suite completed every listed item in both arms with no
OOM. A separate LongCodeQA 1M pilot uses no-KV-cache full-prompt forwards on
identical untruncated prompts.

All runs use PyTorch 2.13.0 (CUDA 13.0), transformers 5.15.1, and the Triton
release bundled with that PyTorch build, together with the corrected R1
tokenizer; results produced before the tokenizer fix are excluded.

\paragraph{Developmental encoder protocol.}
Earlier work used an isolated attention-kernel benchmark (CUDA, 10 timing
trials, 128--32{,}768 tokens) and BART fine-tuning on IMDb at fixed lengths
(256--65{,}536 tokens), recording F1, accuracy, training time, peak memory,
and inference latency. Length and metric coverage vary by architecture;
comparisons are made only at matched lengths. Those runs remain single-seed
estimates and are treated as developmental evidence.

\subsection{Metrics}
\label{sec:metrics}

Table~\ref{tab:metrics} defines the metrics used in this report. For the
current R1 claims, A--D accuracy, prefill time, decode ms/token, and peak
allocated memory are primary. F1 and related IMDb / kernel metrics apply to
the developmental encoder experiments in Section~\ref{sec:developmental}.

\begin{table}[ht]
  \centering
  \small
  \caption{Metrics used in this report.}
  \label{tab:metrics}
  \begin{tabularx}{\textwidth}{@{}lXl@{}}
    \toprule
    Metric & Interpretation & Unit \\
    \midrule
    A--D accuracy & First-token multiple-choice correctness (R1 suites;
      chance $=25\%$) & \% \\
    Prefill time & Full-prompt forward latency & s \\
    Prefill ratio & Mean dense prefill / mean Bigger Bird prefill
      ($>1$ favors BB) & $\times$ \\
    Median paired ratio & Median over examples of dense / BB prefill time;
      robust to a few very long prompts & $\times$ \\
    Decode speed & CUDA-graph probe 
    
    
    & ms/token \\
    Peak memory & Max CUDA allocated over prefill (+ short decode) & GB \\
    F1 & IMDb sentiment quality (developmental) & 0--1 \\
    Attention forward time & Isolated kernel runtime (developmental) & ms \\
    Training time & Wall-clock training time (developmental) & s \\
    OOM boundary & Longest completed context without OOM & tokens \\
    \bottomrule
  \end{tabularx}
\end{table}

Accuracy differences between the paired arms are tested with an exact
McNemar test on per-example correctness. A non-significant result is not
evidence of equivalence, so we report the paired difference alongside the
$p$-value.

\section{Results}
\label{sec:results}
\label{sec:r1-benchmark}

This section reports the primary evidence for Bigger Bird v3: matched
long-context evaluation on R1-Distill-Llama-8B, then the synthetic screens and
ablations that connect earlier prototypes to the final design. Developmental
encoder-scale IMDb / kernel experiments follow in
Section~\ref{sec:developmental} as the rationale and scaling scaffold behind
that path.

\subsection{Matched R1 long-context benchmarks}

The tables below summarize the matched evaluation of Bigger Bird v3 (BB)
against dense attention using BF16 R1-Distill-Llama-8B on a single H100 80 GB GPU (updated November 2, 2026). Both methods were evaluated on 1,178 paired examples across LongCodeQA (repository comprehension), Code.Debug (bug localization), and LongBench v2 (long-context reasoning). Accuracy was measured using first-token multiple-choice predictions (A--D), while prefill latency was measured through full-prompt forwards and decode speed through a separate 16-step CUDA-graph probe. 

\begin{table}[ht]
  \centering
  \small
  \caption{Matched R1 benchmark: dense vs.\ Bigger Bird v3 (H100 80~GB, BF16,
  130K prompt cap). Prefill ratio is mean dense time / mean BB time; above
  $1\times$ favors Bigger Bird.}
  \label{tab:r1-benchmark}
  \begin{tabular}{@{}lrrrr@{}}
    \toprule
    Benchmark & $n$ & Accuracy dense / BB & Prefill dense / BB (s) & Prefill ratio \\
    \midrule
    LongCodeQA 32K  & 113 & 67.3\% / 67.3\% & 0.79 / 1.17 & $0.67\times$ \\
    LongCodeQA 64K  & 76  & 63.2\% / 63.2\% & 2.55 / 3.17 & $0.80\times$ \\
    LongCodeQA 128K & 92  & 52.2\% / 51.1\% & 6.07 / 5.59 & $1.09\times$ \\
    Code.Debug      & 394 & 27.7\% / 26.6\% & 8.75 / 6.89 & $1.27\times$ \\
    LongBench v2    & 503 & 26.8\% / 26.0\% & 6.70 / 5.31 & $1.26\times$ \\
    \bottomrule
  \end{tabular}
\end{table}

The LongCodeQA 128k subset averaged approximately 88K tokens and achieved a median paired prefill ratio of $0.99\times$. Mean decode latency was nearly identical between dense attention and Bigger Bird, ranging from 10.42 to 15.33 ms/token from dense and 10.41 to 15.34 ms/token for Bigger Bird across the five benchmark subsets. Pooled accuracy was 35.3\% for dense attention (416/1178) and 34.6\% for Bigger Bird (407/1178). The paired McNemar test ($p \geq 0.683$) did not establish statistical equivalence. Prompts exceeding 130K were middle-truncated in Code.Debug (136/394) and LongBench~v2 (186/503). These accuracy measurements reflect first-token answer selection rather than successful code patching or repair. 

\begin{table}[ht]
  \centering
  \small
  \caption{KV eviction arm (\texttt{bb\_evict64k}: sink 64 + window 16{,}384 +
  up to 65{,}536 extra selected middle positions). Eviction is applied after
  prefill logits, so first-token QA matches BB v3 by construction.}
  \label{tab:r1-evict}
  \begin{tabular}{@{}lrrr@{}}
    \toprule
    Benchmark & BB / +evict accuracy & Decode BB / +evict (ms/tok) & Peak BB / +evict (GB, mean) \\
    \midrule
    LongCodeQA 32K  & 67.3\% / 67.3\% & 10.41 / 10.53 & 20.9 / 22.9 \\
    LongCodeQA 64K  & 63.2\% / 63.2\% & 12.08 / 12.20 & 25.7 / 30.8 \\
    LongCodeQA 128K & 51.1\% / 51.1\% & 14.08 / 13.71 & 32.3 / 38.4 \\
    Code.Debug      & 26.6\% / 26.6\% & 15.34 / 14.11 & 35.7 / 39.3 \\
    LongBench v2    & 26.0\% / 26.0\% & 13.77 / 12.74 & 31.9 / 34.1 \\
    \bottomrule
  \end{tabular}
\end{table}

\subsection{One-million-token evaluations}

\begin{table}[ht]
  \centering
  \small
  \caption{One-million-token paired evaluations. LongCodeQA: same H100, BF16,
  identical untruncated prompts; paired $p=0.625$. RULER: dense on B200
  (CPU-offloaded KV) vs.\ BB on H100, so no fair prefill ratio.}
  \label{tab:1m-evals}
  \begin{tabular}{@{}lrrrrr@{}}
    \toprule
    Evaluation & $n$ & Tokens & Dense & BB & Prefill (dense $\rightarrow$ BB) \\
    \midrule
    LongCodeQA 1M & 18/50 & 542K--878K & 7/18 (38.9\%) & 9/18 (50.0\%)
      & 470.38 $\rightarrow$ 45.97~s (\textbf{10.23$\times$}) \\
    Synthetic RULER & 18 & 971K--980K & 10/18 (55.6\%) & 9/18 (50.0\%)
      & --- \\
    \bottomrule
  \end{tabular}
\end{table}

\begin{table}[ht]
  \centering
  \small
  \caption{Exp~21 held-out fact retention near 1M tokens (25 prompts,
  4 facts each). No dense BF16 full-cache quality baseline; decode values
  are estimates.}
  \label{tab:1m-exp21}
  \begin{tabular}{@{}lrr@{}}
    \toprule
    Cache setting & Fact recall & Median prefill \\
    \midrule
    Eviction keep 16K  & 62\% & 63~s \\
    Eviction keep 64K  & 85\% & 63~s \\
    Eviction keep 128K & 86\% & 63~s \\
    4-bit full cache   & 91\% & 72--78~s \\
    \bottomrule
  \end{tabular}
\end{table}

% The LongCodeQA test measures first-answer logits from full-prompt forwards without KV cache, rather than decoding or cache retention. Out of 50 prompts, 18 have been evaluated, yielding a median paired prefill speedup of $9.48\times$; the remaining 32 have yet to be tested. 


No KV caching is used in the LongCodeQA experiment. Hence, we only measure first answer logit values. We do not measure 32 prompts. The median ratio of pairs in measured prompts is $9.48\times$.


\subsection{Earlier project screening experiments}

% Table~\ref{tab:screening} records an early prototype screen: byte budget
% 65{,}536, 30 examples, depth 0.5, R1 on an H100 3g.40GB MIG, no KV cache.
% The scorer extracts the first number and checks its last digit. Bigger Bird
% v3 is \emph{not} in this batch; it was developed afterward under a separate
% Modal NIAH protocol (Tables~\ref{tab:v3-niah}--\ref{tab:v3-ablate}).



Table~\ref{tab:screening} shows one of our early screens: byte budget
65{,}536, 30 samples, depth 0.5, R1 on H100 3g.40GB MIG, no KV cache.
The scorer takes the first number from the string and looks at the final digit. Bigger Bird v3 \emph{is not} included here since it was developed later under a different Modal NIAH process
(Tables~\ref{tab:v3-niah}--\ref{tab:v3-ablate}).



\begin{table}[ht]
  \centering
  \small
  \caption{Pre-v3 prototype screening on synthetic needle retrieval
  (30 examples). Relative $=$ last-digit accuracy / dense 30/30. Rows marked
  ``each'' apply the listed score to every named experiment. Development
  record, not a fair ranking of published methods.}
  \label{tab:screening}
  \begin{tabular}{@{}>{\raggedright\arraybackslash}p{0.28\linewidth}
                >{\raggedright\arraybackslash}p{0.42\linewidth}rr@{}}
    \toprule
    Screening result & Experiments & Score & Rel. \\
    \midrule
    Dense reference
      & Exp~0
      & 30/30 & 100\% \\
    Content selection / routing
      & Exp~1 Top-k; 10 GQA; 13 DynamicContext; 15 earlier Bigger Bird (\textbf{each})
      & 27/30 & 90\% \\
    Coarse-to-fine
      & Exp~17
      & 22/30 & 73\% \\
    Bigger Bird Flash
      & Exp~19
      & 16/30 & 53\% \\
    Zero-score prototypes
      & Exp~2, 3, 4, 7, 8, 11, 12, 18 (\textbf{each})
      & 0/30 & 0\% \\
    Not selected here
      & Exp~5, 6, 9, 14, 16
      & --- & --- \\
    \bottomrule
  \end{tabular}
\end{table}

\begin{table}[ht]
  \centering
  \small
  \caption{Bigger Bird v3 synthetic NIAH confirmation (Modal; exact-match
  scoring; depth 0.5 unless noted). Final config:
  \texttt{sink=64}, \texttt{window=16384}, \texttt{vertical=512},
  \texttt{exact\_tail=256}. Prompts 0--5 at $\sim$127K were used in tuning;
  held-out is prompts 6--29.}
  \label{tab:v3-niah}
  \begin{tabular}{@{}lrrrrr@{}}
    \toprule
    Context & Arm & Exact (30) & Held-out (24) & Prefill (s) & Ratio \\
    \midrule
    $\sim$65K  & Dense & 30/30 & 24/24 & 5.02 & --- \\
    $\sim$65K  & BB v3 & 30/30 & 24/24 & 4.41 & $1.14\times$ \\
    $\sim$128K & Dense & 28/30 & 22/24 & 16.25 & --- \\
    $\sim$128K & BB v3 & 30/30 & 24/24 & 8.61 & $1.89\times$ \\
    $\sim$128K (depth 0.1) & Dense & 28/30 & 23/24 & 16.39 & --- \\
    $\sim$128K (depth 0.1) & BB v3 & 30/30 & 24/24 & 8.56 & $1.91\times$ \\
    \bottomrule
  \end{tabular}
\end{table}

\begin{table}[ht]
  \centering
  \small
  \caption{Why Exp~19 underperformed relative to v3: matched tuning-set
  prompts at $\sim$127K (depth 0.5, $n{=}6$). Exact-tail ablation uses a 4K
  window; final v3 uses the 16K window.}
  \label{tab:v3-ablate}
  \begin{tabular}{@{}lrr@{}}
    \toprule
    Configuration & Exact & Prefill (s) \\
    \midrule
    Dense & 6/6 & 16.33 \\
    Exp~19 Bigger Bird Flash & 0/6 & 4.33 \\
    Exact tail $=0$ (4K window) & 2/6 & 4.95 \\
    Exact tail $=64$ (4K window) & 4/6 & 5.45 \\
    Exact tail $=256$ (4K window) & 6/6 & 6.09 \\
    BB v3 final (16K window) & 6/6 & 8.50 \\
    \bottomrule
  \end{tabular}
\end{table}

Table~\ref{tab:v3-niah} shows that the final v3 config recovers dense-level
needle retrieval while nearly doubling prefill speed at $\sim$128K.
Table~\ref{tab:v3-ablate} links Exp~19's weak screen to missing coverage:
without a sufficient exact tail, the same development prompts fail even when
prefill is fast. These synthetic results motivate the design; the primary
real-benchmark claim remains Section~\ref{sec:r1-benchmark}.


\section{Developmental Encoder Experiments}
\label{sec:developmental}

Before moving to decoder models, we explored sparse attention on
encoder-scale models (BART and BigBird-RoBERTa) fine-tuned for IMDb
sentiment classification. These experiments are smaller and less controlled
than those in Section~\ref{sec:results}, but several design choices in Bigger
Bird~v3 came directly from them, and we report them for that reason.

Two limitations apply throughout this section. First, most IMDb reviews are a
few hundred tokens long, so at 8K tokens and beyond the inputs consist largely
of padding. These runs therefore measure cost and training stability at
length, not the use of long-range evidence. Second, each configuration was
run once, with training budgets that differed across runs, so F1 differences
of a few points fall within the variation we observed between neighbouring
lengths.

Section~\ref{sec:dev-outcomes} summarizes each prototype and what it showed.
The following subsections report kernel and task-level efficiency, and
Section~\ref{sec:bigbird-family-results} compares selection rules within the
BigBird family.

\subsection{Per-architecture outcomes}
\label{sec:dev-outcomes}

Table~\ref{tab:dev-summary} lists every encoder-scale prototype. Most of the
results fall into four groups, and each group informed a different part of
the decoder design.

\paragraph{Learned routers were fast but unstable.}
Several prototypes added a small learned component to decide where each query
should look: a low-rank key scorer (Experiment~1), a gate for global tokens
(Experiment~3), a block router (Experiment~4), and combinations of these
(Experiments~6 and~10). Their kernels were the fastest we measured, about
17$\times$ dense attention at 16K tokens, but training often collapsed to
majority-class prediction. Experiment~1 fell to F1 0.519 at 2K tokens and to
zero at 32K, and Experiment~6 collapsed at every length tested. The decoder
screens in Section~\ref{sec:results} showed a similar pattern, with untrained
selection components failing retrieval. For this reason, Bigger Bird~v3 does
not add a learned router; it selects distant keys using the attention scores
of the pretrained model itself.

\paragraph{Kernel speedups did not carry over end to end.}
Layer-Adaptive Sparsity (Experiment~7) was the most stable sparse method,
within 0.003 F1 of dense attention at 1K tokens, and its kernel ran
17.5$\times$ faster than dense at 16K. However, inference at that length was
only about 1.5$\times$ faster (169~ms versus 254~ms per sequence), and at 4K
tokens a 97.7\% reduction in query--key comparisons produced a 12\% latency
gain. The rest of the network, together with the cost of gathering selected
keys, absorbed most of the saving. As a result, all speed results for v3 are
reported as full-prompt prefill times on the complete model rather than as
isolated kernel times.

\paragraph{Shortening the sequence saved the most resources.}
Token Dropping (Experiment~8) discards 30\% of tokens after three dense
layers. At 32K tokens it ran 3.47$\times$ faster than dense attention with
82.5\% less peak memory, at a cost of 0.048 F1, and it was the only method to
complete a 65K-token run. Dynamic Context (Experiment~13) goes further and
keeps a fixed budget of 4{,}096 tokens. It reached the highest IMDb score (F1
0.942 at 1K) and retrieved every needle on RULER at 4K, but its F1 fell to
0.03 or below once inputs exceeded that budget. This suggests that a fixed
token budget is effective only while the relevant evidence remains within it,
which is why Bigger Bird~v3 keeps every token in the KV cache and gives the
final prompt queries exact attention.

\paragraph{The first Bigger Bird preserved quality but was costly.}
Unified Bigger Bird (Experiment~5) combined MMR local selection, learned
globals, and teleports on BART. It matched dense quality at 768 tokens (F1
0.926 versus 0.931), but it trained 12.8$\times$ more slowly, ran slower than
dense in the 1K--2K kernel benchmark, and ran out of memory at 4K. Although
the selection ideas preserved quality, the implementation was too costly to
scale. The later Triton prototypes (Experiment~19 and v3) moved selection
into a fused kernel that does not build the full attention matrix.

\begin{table}[ht]
  \centering
  \small
  \caption{Encoder-scale prototypes on IMDb (single runs). Kernel speedups
  are isolated attention time relative to dense. ``Collapse'' means
  majority-class prediction (F1 = 0). Per-length results are in
  Table~\ref{tab:app-imdb}.}
  \label{tab:dev-summary}
  \begin{tabularx}{\textwidth}{@{}l>{\raggedright\arraybackslash}X@{}}
    \toprule
    Experiment & Outcome \\
    \midrule
    0 Dense & Best F1 0.931 (768 tokens); out of memory at 65K \\
    \midrule
    1 Low-rank top-$k$ & Kernel 17.4$\times$ at 16K; F1 0.519 at 2K, collapse at 32K \\
    3 Dynamic globals & Strong at 512 tokens; collapses at several longer lengths \\
    4 Block routing (PBS) & Kernel faster than dense from 1K; several collapses \\
    6 Top-$k$ + block routing & Collapse at every length from 512 to 4K \\
    10 GQA + top-$k$ & Kernel about 17$\times$ at 16K; mixed useful and collapsed runs \\
    \midrule
    7 Layer-adaptive & Within 0.003 F1 of dense at 1K; kernel 17.5$\times$ at 16K \\
    8 Token dropping & 3.47$\times$ faster, 82.5\% less memory at 32K; only 65K run \\
    13 Dynamic context & F1 0.942 at 1K; about 10$\times$ training time; fails retrieval past 4K \\
    14 Token drop + top-$k$ & Competitive on LRA Text to 4K; out of memory beyond \\
    \midrule
    2 Lightning hybrid & Kernel slower than dense; several collapses \\
    9 Attention speculation & Out of memory at 4K \\
    11 NSA & One strong short run; implementation unfinished \\
    12 S2-HHST & Moderate short-context F1; implementation unfinished \\
    \midrule
    5 Unified Bigger Bird & F1 0.926 versus 0.931 at 768; 12.8$\times$ training; out of memory at 4K \\
    15 Proper Bigger Bird & No standalone run (see Section~\ref{sec:bigbird-family-results}) \\
    \bottomrule
  \end{tabularx}
\end{table}

\subsection{Highlighted models and early scaling plots}


The developmental experiments provide two complementary views of efficiency: isolated attention-kernel execution time and task-level performance. We highlight the Dense Baseline (Experiment 0), Unified Bigger Bird (Experiment 5), Dynamic Context (Experiment 13), and Token Dropping (Experiment 8) in the task-level comparisons. These represent the reference implementation, the initial Bigger Bird architecture, and two alternative approaches to reducing attention cost.

For kernel-level comparisons, Layer-Adaptive Sparsity (Experiment 7) replaces Dynamic Context, for which no isolated kernel measurements were recorded. Additional sparse attention methods are included in the numerical tables to provide a broader comparison of scaling behaviour.

These experiments reveal several trade-offs. Layer-Adaptive Sparsity achieves substantial kernel-level acceleration, while Token Dropping provides stronger end-to-end memory and latency reductions at longer sequence lengths. Dynamic Context achieves high short-context classification accuracy but introduces a fixed token budget. In contrast, Unified Bigger Bird preserves classification quality while incurring substantial computational overhead.

Because the experiments use single runs with differing training budgets, and longer IMDb inputs contain substantial padding, the results should be interpreted as exploratory evidence of computational behaviour rather than definitive comparisons of model quality.

\subsection{Log-Log Attention Forward Time}
\label{sec:loglog}

\begin{figure}[ht]
  \centering
  \includegraphics[width=\linewidth]{sel_loglog}
  \caption{Log--log attention-kernel time for four experiments: Dense
  Baseline (0), Bigger Bird Unified (5), Layer-Adaptive Sparsity (7), and
  Token Dropping (8). Layer-Adaptive, Token Dropping, and the dense baseline
  extend to 32{,}768 tokens; the dense baseline OOMs at that length.
  Dynamic Context (13) is absent because it was not included in the kernel
  benchmark; it appears in the efficiency figures (Figure~\ref{fig:efficiency}).
  Kernel times for all other experiments appear in Table~\ref{tab:app-kernel}.}
  \label{fig:loglog}
\end{figure}


Figure~\ref{fig:loglog} illustrates attention-kernel execution time as a function of sequence length on logarithmic axes. Table~\ref{tab:loglog} provides additional measurements for the Dense Baseline, DeepSeek Top-k, Layer-Adaptive Sparsity, and GQA Sparse Attention. The figure and table include different subsets of the developmental architectures.

At short sequence lengths, the overhead of sparse attention can outweigh the benefits of reducing query-key interactions. At 128 tokens, dense attention requires 0.386 ms, compared with 0.463--0.528 ms for the three sparse methods in Table~\ref{tab:loglog}. As sequence length increases, sparse attention becomes more advantageous. At 4096 tokens, the sparse methods achieve approximately 3.9--4.6 times lower kernel execution time than dense attention.

At 16,384 tokens, dense attention requires 95.415 ms, while the three sparse methods complete in approximately 5.5 ms, corresponding to speedups of 17.1--17.5 times. All three sparse methods also complete the 32,768-token kernel benchmark in approximately 10.5--10.7 ms, whereas the dense implementation encounters an out-of-memory failure.

% numbers look a bit strange for figure 1/tabel 12, time decreases between 4k and 8k tokens; needs to be re-tested. Should be whatever for the manuscript. There is a note about this discontinuity but it more-so just states the way it was measured was inconsistent: "This inversion is inconsistent with monotonic quadratic growth and likely reflects a difference between the original and extended timing sweeps rather than a genuine speedup". I removed that part and the dagger in the table below as it does not provide anything. 

% TLDR: Basically we have no good explanation for this and should be re-tested. 


\begin{table}[ht]
  \centering
  \small
  \caption{Attention-forward time (ms) across nearly the full tested range,
  log-scaled sequence length. OOM denotes an out-of-memory failure.}
  \label{tab:loglog}
  \begin{tabular}{@{}rrrrr@{}}
    \toprule
    Seq.\ length & Dense (Exp 0) & DeepSeek Top-k (1) & Layer-Adaptive (7) & GQA Sparse (10) \\
    \midrule
    128    & 0.386 & 0.463 & 0.528 & 0.493 \\
    256    & 0.632 & 0.487 & 0.487 & 0.503 \\
    512    & 1.396 & 0.858 & 0.857 & 0.883 \\
    1024   & 3.537 & 1.567 & 1.567 & 1.595 \\
    2048   & 9.007 & 2.889 & 2.897 & 2.961 \\
    4096   & 27.092 & 5.834 & 6.707 & 6.925 \\
    %8192   & 26.263$^\dagger$ & 2.935 & 2.953 & 3.007 \\
    8192   & 26.263 & 2.935 & 2.953 & 3.007 \\
    16384  & 95.415 & 5.485 & 5.464 & 5.589 \\
    32768  & OOM & 10.522 & 10.504 & 10.739 \\
    \bottomrule
  \end{tabular}

  %\vspace{4pt}
  %\raggedright\footnotesize $^\dagger$Lower than the 4096-token value for the
  %same method; see discussion of the extended-sweep discontinuity above.
  % ^ ? 
\end{table}

\subsection{Linear Scale on Short and Long Sequences} 

\begin{figure}[ht]
  \centering
  \begin{subfigure}[b]{0.48\linewidth}
    \includegraphics[width=\linewidth]{sel_linear_short}
    \caption{Short sequences (128--512 tokens).}
    \label{fig:linshort}
  \end{subfigure}
  \hfill
  \begin{subfigure}[b]{0.48\linewidth}
    \includegraphics[width=\linewidth]{sel_linear_long}
    \caption{Long sequences (1024--4096 tokens).}
    \label{fig:linlong}
  \end{subfigure}
  \caption{Attention-kernel time on a linear scale for Dense Baseline (0),
  Bigger Bird Unified (5), Layer-Adaptive Sparsity (7), and Token Dropping (8).
  Dynamic Context (13) has no kernel benchmark and is omitted from these panels;
  it appears in Figure~\ref{fig:efficiency}. Kernel times for all other
  experiments appear in Table~\ref{tab:app-kernel}.}
  \label{fig:linear}
\end{figure}

Figures~\ref{fig:linshort} and~\ref{fig:linlong} present kernel execution times on linear axes, separately highlighting short and long sequence lengths. Tables~\ref{tab:short} and~\ref{tab:long} provide the corresponding measurements for the broader set of developmental architectures.

\begin{table}[ht]
  \centering
  \small
  \caption{Attention-forward time (ms), short sequences (128--512 tokens),
  linear scale.}
  \label{tab:short}
  \begin{tabular}{@{}lrrr@{}}
    \toprule
    Method & 128 & 256 & 512 \\
    \midrule
    Dense baseline (0)        & 0.386 & 0.632 & 1.396 \\
    DeepSeek Top-k (1)        & 0.463 & 0.487 & 0.858 \\
    Lightning Hybrid (2)      & 1.157 & 2.209 & 4.277 \\
    Dynamic Globals (3)       & 0.873 & 1.432 & 2.642 \\
    PBS Attn (4)              & 0.307 & 0.580 & 0.915 \\
    Bigger Bird Unified (5)   & 1.078 & 1.872 & 3.499 \\
    DeepSeek+PBS (6)          & 0.809 & 0.783 & 0.974 \\
    Layer-Adaptive (7)        & 0.528 & 0.487 & 0.857 \\
    Token Drop (8)            & 0.303 & 0.499 & 1.105 \\
    Attn Speculation (9)      & 0.745 & 1.294 & 2.433 \\
    GQA+Sparse (10)           & 0.493 & 0.503 & 0.883 \\
    \bottomrule
  \end{tabular}
\end{table}

\begin{table}[ht]
  \centering
  \small
  \caption{Attention-forward time (ms), long sequences (1024--4096 tokens),
  linear scale.}
  \label{tab:long}
  \begin{tabular}{@{}lrrr@{}}
    \toprule
    Method & 1024 & 2048 & 4096 \\
    \midrule
    Dense baseline (0)        & 3.537 & 9.007  & 27.092 \\
    DeepSeek Top-k (1)        & 1.567 & 2.889  & 5.834 \\
    Lightning Hybrid (2)      & 8.783 & 17.406 & 34.411 \\
    Dynamic Globals (3)       & 5.068 & 9.940  & 19.447 \\
    PBS Attn (4)              & 1.663 & 3.068  & 6.162 \\
    Bigger Bird Unified (5)   & 6.798 & 13.455 & 26.502 \\
    DeepSeek+PBS (6)          & 1.678 & 3.014  & 6.147 \\
    Layer-Adaptive (7)        & 1.567 & 2.897  & 6.707 \\
    Token Drop (8)            & 2.823 & 8.315  & 27.085 \\
    Attn Speculation (9)      & 4.722 & 9.451  & 18.575 \\
    GQA+Sparse (10)           & 1.595 & 2.961  & 6.925 \\
    \bottomrule
  \end{tabular}
\end{table}


At short sequence lengths, sparse attention does not provide a universal performance advantage. At 128 tokens, DeepSeek Top-k and Layer-Adaptive Sparsity are slower than dense attention, whereas PBS Attention and Token Dropping are slightly faster. By 256 tokens, DeepSeek Top-k, Layer-Adaptive Sparsity, and GQA Sparse Attention outperform the dense baseline. Other architectures, including Lightning Hybrid and Unified Bigger Bird, retain substantial overhead.

The advantage of efficient sparse kernels becomes more pronounced at longer sequence lengths. Between 1024 and 4096 tokens, dense attention execution time increases from 3.537 ms to 27.092 ms, a factor of approximately 7.7. Over the same range, DeepSeek Top-k increases from 1.567 ms to 5.834 ms, a factor of approximately 3.7. This difference results in a growing relative advantage for the sparse implementation.

Nevertheless, not all sparse architectures achieve improved kernel performance. Lightning Hybrid remains slower than dense attention throughout the 1024--4096-token range. Unified Bigger Bird is also slower at 1024 and 2048 tokens, reaching only a marginal advantage at 4096 tokens (26.502 ms versus 27.092 ms).

Token Dropping similarly approaches dense kernel execution time at 4096 tokens
(27.085 ms versus 27.092 ms). Its isolated attention measurements therefore understate the benefits observed at the full-model level, where dropping tokens reduces subsequent-layer computation and memory requirements.

Together, these results demonstrate that reducing attention operations does not automatically translate into proportional execution-time improvements. Routing overhead, memory access patterns, and the fraction of computation affected by sparsification all influence realized performance.

\subsection{Efficiency Metrics: F1, Latency, Memory, and Training Time} % @thomas In general, the tables and graphs here are a bit strange because they include and exclude various experiments seemingly at random
\label{sec:efficiency}

\begin{figure}[ht]
  \centering
  \begin{subfigure}[b]{0.49\linewidth}
    \includegraphics[width=\linewidth]{sel_f1}
    \caption{IMDb F1 vs sequence length.}
    \label{fig:f1seq}
  \end{subfigure}
  \hfill
  \begin{subfigure}[b]{0.49\linewidth}
    \includegraphics[width=\linewidth]{sel_latency}
    \caption{Inference latency vs sequence length.}
    \label{fig:latseq}
  \end{subfigure}

  \medskip
  \begin{subfigure}[b]{0.49\linewidth}
    \includegraphics[width=\linewidth]{sel_memory}
    \caption{Peak device memory vs sequence length.}
    \label{fig:memseq}
  \end{subfigure}
  \hfill
  \begin{subfigure}[b]{0.49\linewidth}
    \includegraphics[width=\linewidth]{sel_traintime}
    \caption{Training time vs sequence length.}
    \label{fig:trainseq}
  \end{subfigure}

  \caption{IMDb task-level efficiency for the four highlighted experiments:
  Dense Baseline (0), Bigger Bird Unified (5), and the two top-performing
  alternatives by overall score: Dynamic Context (13) and Token Dropping (8).
  Overall score is the mean of six normalized axes --- F1, accuracy, inference
  speed, memory efficiency, training efficiency, and sparsity --- averaged
  across all recorded sequence lengths. Full results for all remaining
  experiments appear in Table~\ref{tab:app-imdb}.}
  \label{fig:efficiency}
\end{figure}

Figure~\ref{fig:efficiency} compares the IMDb task-level F1 scores, latency, memory, and training time for the dense baseline, Bigger Bird, and two of the best performing sparse attention models. Tables~\ref{tab:eff-baseline}--\ref{tab:eff-dynamiccontext} include the DeepSeek Top-k model because a broad range of tests were performed for it; all
other results appear in Table~\ref{tab:app-imdb}. An em-dash denotes a missing
metric, and OOM denotes an out-of-memory failure.

\begin{table}[ht]
  \centering
  \small
  \caption{Dense baseline (Experiment 0): IMDb efficiency by sequence length.}
  \label{tab:eff-baseline}
  \begin{tabular}{@{}rrrrr@{}}
    \toprule
    Seq.\ length & F1 & Train time (s) & Peak memory (MB) & Latency (ms) \\
    \midrule
    256    & 0.878 & 157.3   & ---      & ---   \\
    512    & 0.871 & 56.3    & 2707.4   & 13.68 \\
    768    & 0.931 & 216.7   & ---      & ---   \\
    1024   & 0.894 & 57.3    & 2725.4   & 13.18 \\
    2048   & 0.900 & 66.6    & 3061.0   & 13.21 \\
    4096   & 0.909 & 135.6   & 4446.4   & 24.47 \\
    8192   & 0.871 & 441.4   & 3408.8   & 75.72 \\
    16384  & 0.917 & 1484.9  & 6793.8   & 253.99 \\
    32768  & 0.840 & 5485.6  & 20{,}577.4 & 924.88 \\
    65536  & OOM   & OOM     & OOM      & OOM   \\
    \bottomrule
  \end{tabular}
\end{table}

\begin{table}[ht]
  \centering
  \small
  \caption{DeepSeek Top-k (Experiment 1): IMDb efficiency by sequence length.}
  \label{tab:eff-deepseek}
  \begin{tabular}{@{}rrrrr@{}}
    \toprule
    Seq.\ length & F1 & Train time (s) & Peak memory (MB) & Latency (ms) \\
    \midrule
    256    & 0.917 & 187.2  & ---      & ---   \\
    512    & 0.860 & 60.8   & 2701.4   & 16.07 \\
    768    & 0.926 & 609.8  & ---      & ---   \\
    1024   & 0.820 & 61.6   & 2726.6   & 15.00 \\
    2048   & 0.519 & 62.5   & 3077.6   & 15.08 \\
    4096   & 0.821 & 102.6  & 4385.7   & 20.40 \\
    8192   & 0.600 & 295.6  & 3390.2   & 53.94 \\
    16384  & 0.812 & 907.0  & 6791.6   & 168.21 \\
    32768  & 0.000 & 3170.7 & 20{,}577.1 & 581.10 \\
    65536  & OOM   & OOM    & OOM      & OOM   \\
    \bottomrule
  \end{tabular}
\end{table}

\begin{table}[ht]
  \centering
  \small
  \caption{Layer-Adaptive Sparsity (Experiment 7): IMDb efficiency by
  sequence length.}
  \label{tab:eff-layeradaptive}
  \begin{tabular}{@{}rrrrr@{}}
    \toprule
    Seq.\ length & F1 & Train time (s) & Peak memory (MB) & Latency (ms) \\
    \midrule
    512    & 0.800 & 118.1  & 2696.0   & 19.58 \\
    1024   & 0.891 & 84.5   & 2724.5   & 15.95 \\
    2048   & 0.830 & 81.3   & 2781.5   & 15.59 \\
    4096   & 0.864 & 124.6  & 2817.0   & 21.88 \\
    8192   & 0.826 & 298.1  & 3406.6   & 54.63 \\
    16384  & 0.842 & 911.5  & 6792.8   & 169.32 \\
    32768  & 0.778 & 3178.2 & 20{,}575.9 & 582.94 \\
    65536  & OOM   & OOM    & OOM      & OOM   \\
    \bottomrule
  \end{tabular}
\end{table}

\begin{table}[ht]
  \centering
  \small
  \caption{Token Dropping (Experiment 8): IMDb efficiency by sequence length.}
  \label{tab:eff-tokendrop}
  \begin{tabular}{@{}rrrrr@{}}
    \toprule
    Seq.\ length & F1 & Train time (s) & Peak memory (MB) & Latency (ms) \\
    \midrule
    1024   & 0.286 & 37.5   & 1833.6   & 5.41 \\
    2048   & 0.776 & 41.3   & 1857.3   & 5.21 \\
    4096   & 0.808 & 64.8   & 1894.2   & 8.33 \\
    8192   & 0.837 & 159.9  & 1968.2   & 23.09 \\
    16384  & 0.851 & 503.3  & 2494.4   & 74.48 \\
    32768  & 0.792 & 1837.5 & 3603.6   & 266.92 \\
    65536  & 0.780 & 7112.3 & 5815.3   & 1019.87 \\
    \bottomrule
  \end{tabular}
\end{table}

\FloatBarrier
\begin{table}[H]
  \centering
  \small
  \caption{Dynamic Context (Experiment 13): IMDb efficiency by sequence
  length.}
  \label{tab:eff-dynamiccontext}
  \begin{tabular}{@{}rrrrr@{}}
    \toprule
    Seq.\ length & F1 & Train time (s) & Peak memory (MB) & Latency (ms) \\
    \midrule
    256    & 0.000 & 26.9  & 1826.8 & 5.31 \\
    512    & 0.909 & 65.5  & 1862.3 & 5.20 \\
    768    & 0.919 & 107.4 & 2311.6 & 5.30 \\
    1024   & 0.942 & 566.4 & 3457.7 & 5.35 \\
    4096   & 0.788 & 74.9  & 1898.6 & 9.85 \\
    \bottomrule
  \end{tabular}
\end{table}

Several patterns emerge across Tables~\ref{tab:eff-baseline}--\ref{tab:eff-dynamiccontext}.
Inference latency separates the methods most clearly: Dynamic Context holds latency in the 5.2--9.9~ms range up to a context length of up to 4096, roughly 2.5--2.6$\times$ lower than dense attention at the same length, because its fixed post-drop token budget bounds the amount of computation regardless of input length. Token Dropping shows the same qualitative pattern for inference latency at longer
lengths. In addition, starting at around 8192 tokens, the peak memory usage for Token Dropping grows substantially more slowly than that of dense attention (Fig.~\ref{fig:efficiency}). 

DeepSeek Top-k and Layer-Adaptive Sparsity track dense attention's memory footprint closely at short-to-medium lengths (within a few percent at 512--2048 tokens) because their savings are concentrated in query-key comparison count rather than in the number of activations % TODO@thomas activation functions?
stored, and their latency advantage over dense attention only becomes large
at 8192 tokens and beyond (53.9--54.6~ms versus 75.7~ms).
% @thomas Better to make a graph for Tables 15-19

Quality (F1) is comparatively more volatile than efficiency: two of the five methods (DeepSeek Top-k and Token
Dropping) each have at least one sequence length where F1 drops sharply below neighboring lengths (DeepSeek Top-k: 0.519 at 2048 tokens versus 0.820 at 1024 and 0.821 at 4096; Token Dropping: 0.286 at 1024 tokens versus 0.776 at 2048).
These dips do not follow a consistent trend with sequence length and are most plausibly attributed to training-run variance (sample count, epoch count, and random seed all differ across individual runs) rather than to a
property of the mechanism itself; they argue for the matched, multi-seed re-run recommended in Section~\ref{sec:limitations}.

Table~\ref{tab:quality-1024} isolates the 1024-token comparison directly, and
Table~\ref{tab:tradeoffs} restates the same data as relative speedups and
memory reductions against the dense baseline at representative lengths.

\begin{table}[H]
  \centering
  \small
  \caption{IMDb quality at sequence length 1024.}
  \label{tab:quality-1024}
  \begin{tabular}{@{}lrrr@{}}
    \toprule
    Method & F1 & $\Delta$F1 vs.\ dense & Accuracy \\
    \midrule
    Dense baseline & 0.894 & 0.000 & 0.900 \\
    Dynamic Context & \textbf{0.942} & +0.048 & \textbf{0.941} \\
    Layer-Adaptive & 0.891 & $-0.003$ & 0.900 \\
    Attention Speculation & 0.841 & $-0.053$ & 0.860 \\
    DeepSeek Top-k & 0.820 & $-0.074$ & 0.820 \\
    Token Drop & 0.286 & $-0.608$ & 0.600 \\
    \bottomrule
  \end{tabular}
\end{table}

\begin{table}[H]
  \centering
  \small
  \caption{Representative quality--efficiency trade-offs relative to the dense
  baseline. Speedups above one are favorable; negative memory reduction denotes
  an increase.}
  \label{tab:tradeoffs}
  \begin{tabular}{@{}rrrcrr@{}}
    \toprule
    Length & Method & $\Delta$F1 & Inference speedup & Memory reduction & Train speedup \\
    \midrule
    1024  & Dynamic Context & +0.048 & 2.46$\times$ & $-26.9$\% & 0.10$\times$ \\
    4096  & Layer-Adaptive & $-0.045$ & 1.12$\times$ & 36.6\% & 1.09$\times$ \\
    8192  & Token Drop & $-0.034$ & 3.28$\times$ & 42.3\% & 2.76$\times$ \\
    16384 & Token Drop & $-0.066$ & 3.41$\times$ & 63.3\% & 2.95$\times$ \\
    32768 & Token Drop & $-0.048$ & 3.47$\times$ & 82.5\% & 2.99$\times$ \\
    \bottomrule
  \end{tabular}
\end{table}
% @thomas Would be nice if we had a table or graph for "comparisons"
Reported query-key comparison counts do not solely predict measured latency. At 1,024 tokens, Dynamic Context reports about 1.85 billion comparisons, compared to 151 million for dense attention, yet its measured latency is lower. At 4,096 tokens, Layer-Adaptive Sparsity reports a reduction in query-key comparisons by 97.7\%, but improves end-to-end latency by only 12\%. This shows that runtime also depends on other factors such as kernel overhead, when tokens are dropped, and what the timing includes. The query-key comparison count metric should be defined alongside these results.

\FloatBarrier
\subsection{BigBird-family developmental comparison}
\label{sec:bigbird-family-results}

The main bet behind Bigger Bird is that long-range links should be chosen
from the content instead of being fixed in advance. This study was our first direct test
of that bet. Starting from the original BigBird \citep{zaheer2020bigbird}, which
mixes local, global, and random links, we built BART-based variants that added
MMR local selection and facility-location globals, then compared random
links against targeted teleports (Section~\ref{exp:bigbird-family}).
Figure~\ref{fig:bigbirdfamily} shows the results at 768 and 8{,}000 tokens.

\begin{figure}[ht]
  \centering
  \includegraphics[width=\linewidth]{fig_bigbird_family}
  \caption{BigBird-family IMDb F1 at 768 tokens (left) and 8{,}000 tokens
  (right).}
  \label{fig:bigbirdfamily}
\end{figure}

At 768 tokens the variants were indistinguishable. Every one stayed within
0.005 F1 of dense attention (Table~\ref{tab:bigbird-family}), so adding
global selection or MMR neither helped nor hurt at this length.

\begin{table}[ht]
  \centering
  \small
  \caption{BigBird-family IMDb results at 768 tokens (single runs).}
  \label{tab:bigbird-family}
  \begin{tabular}{@{}lrrr@{}}
    \toprule
    Method & F1 & Accuracy & Training time (s) \\
    \midrule
    Dense baseline & 0.931 & 0.931 & 216.7 \\
    BigBird & 0.928 & 0.931 & 2766.6 \\
    Bigger Bird: globals & 0.926 & 0.929 & 3055.8 \\
    Bigger Bird: top-$k$ MMR & 0.927 & 0.930 & 2902.7 \\
    Bigger Bird: top-$k$ MMR v2 & 0.929 & 0.932 & 2929.3 \\
    Bigger Bird: top-$k$ MMR + globals & 0.929 & 0.932 & 3184.5 \\
    \bottomrule
  \end{tabular}
\end{table}

At 8{,}000 tokens the variants separated (Table~\ref{tab:bigbird-8000}).
Adding teleports to MMR and globals raised F1 to 0.940, the best in the
family, against 0.925 for the unmodified BigBird. Adding random links instead gave
0.928, slightly below the same configuration without them (0.938). These results are
from single runs, so the gaps are suggestive rather than established, but their
direction matters for this paper as they suggest that links chosen by content
help, and links chosen at random do not.

\begin{table}[ht]
  \centering
  \small
  \caption{BigBird-family IMDb results at 8{,}000 tokens (single runs). No
  dense run was recorded at exactly this length.}
  \label{tab:bigbird-8000}
  \begin{tabular}{@{}lrrr@{}}
    \toprule
    Method & F1 & Accuracy & Training time (s) \\
    \midrule
    BigBird & 0.925 & 0.928 & 21{,}504.5 \\
    Bigger Bird: top-$k$ MMR + globals & 0.938 & 0.940 & 7484.3 \\
    Bigger Bird: top-$k$ MMR + globals + random & 0.928 & 0.930 & 5832.9 \\
    Bigger Bird: top-$k$ MMR + globals + teleports & \textbf{0.940} & \textbf{0.941} & \textbf{4348.7} \\
    \bottomrule
  \end{tabular}
\end{table}

However, these quality gains came at a substantial training cost. Every
variant trained far more slowly than dense attention: 12--14.7$\times$ at 768
tokens, and at 8{,}000 tokens the teleport variant took 4{,}349~s against
441~s for the closest dense run (8{,}192 tokens,
Table~\ref{tab:eff-baseline}). Teleports did cut training time by 4.9$\times$ relative to stock BigBird,
but not enough to close that gap. These results motivated the design of
Bigger Bird~v3, which keeps content-based selection of distant positions but
removes the training cost by applying it at inference time, through a Triton
kernel, to a pretrained decoder.

\FloatBarrier
\section{Discussion}

\subsection{Empirical Synthesis and Design Trajectory}

Our results suggest that Bigger Bird v3 offers a promising trade-off between long-context inference and answer quality. The matched R1 evaluation showed that sparse prefill maintained accuracy close to dense attention while achieving mean speedups of 1.09-1.27x on longer benchmarks. However, sparse attention was slower at shorter context lengths, indicating that routing overhead can outweigh computational savings until sequences become sufficiently long. The significantly larger speedup of 10.23x observed in the near 1M token pilot further suggests that the approach may become more advantageous at extreme context lengths, although this result remain preliminary. 

The developmental IMDb / kernel experiments highlight a common distinction between theoretical sparsity, kernel efficiency, and end-to-end performance. Several early architectures achieved substantial kernel speedups but suffered from unstable task quality or limited overall latency improvements. 

DeepSeek Top-k, Layer-Adaptive, and GQA Sparse showed strong
isolated kernel scaling; Token Dropping and Dynamic Context offered distinct quality--efficiency profiles. Subsequent R1 experiments showed that untrained routing mechanisms and insufficient attention budgets can compromise long-range information retrieval. Hence, these findings motivated Bigger Bird v3's combination of attention sinks, local windows, content-based routing, and an exact attention tail. Overall, the results suggest that effective sparse attention requires balancing computational savings against the preservation of task-relevant information, rather than maximizing sparsity alone. 

\subsection{Complexity and Break-Even Analysis}

For a prompt of length $n$, standard causal self-attention evaluates $\frac{n(n+1)}{2} = \mathcal{O}(n^2)$ query--key inner products per head in each layer~\cite{vaswani2017attention}. Bigger Bird~v3 parameterizes the sparse prefill mask with five fixed constants: $S = 64$ attention sinks, a local causal window of width $W = 16{,}384$, up to $R = 512$ content-routed distant positions per head, a late-prompt routing query set of size $Q_r = 64$, and an exact dense tail of length $T = 256$.

Whenever $n > S + W + T$, the total number of per-head query--key inner
products in a single prefill layer is bounded by four terms:
\begin{align}
C_{\mathrm{BBv3}}(n) &\le
  \underbrace{\frac{(S+W)(S+W+1)}{2}}_{\text{prefix dense corner}}
  + \underbrace{(n - T - S - W)(S + W + R)}_{\text{sparse prefill body}} \nonumber \\
  &\quad + \underbrace{Tn - \frac{T(T-1)}{2}}_{\text{exact tail}}
  + \underbrace{Q_r n - \frac{Q_r(Q_r-1)}{2}}_{\text{routing pass}}.
\label{eq:bb-comparisons}
\end{align}
The first term accounts for queries $i \le S + W$, whose sinks and sliding window already cover their full causal prefix. The second term bounds the sparse body ($S + W < i \le n - T$), where each query attends to at most $S + W + R = 16{,}960$ keys. The third term reflects exact causal attention for the final $T = 256$ prompt queries, and the fourth bounds the cost of scoring all prefix keys against the final $Q_r = 64$ queries to select the routed indices (an overhead of less than $0.4\%$ of the sparse body). Asymptotically, the total comparison count scales linearly as $\mathcal{O}\bigl(n(S + W + R + T + Q_r)\bigr) = \mathcal{O}(n)$ per layer. Collapsing and combining terms, the exact cost bound for $n > 16{,}704$ is $$C_{\mathrm{BBv3}}(n) \le 17{,}280\,n - 1.481 \times 10^8.$$

Comparing $\frac{n(n+1)}{2}$ against $C_{\mathrm{BBv3}}(n)$ yields a theoretical attention-comparison break-even point at $n \approx 18{,}850$ tokens. Because attention represents only a fraction of full-model prefill time and dense SDPA benefits from highly optimized FlashAttention Tensor Core scheduling on GPUs \cite{dao2022flashattention}, the empirical crossover shown in Tables~\ref{tab:r1-benchmark} and \ref{tab:1m-evals} comes much later at around $\sim100{,}000$ context size, where the theoretical attention comparison ratio is $\sim3\times$. This is when the reduction in FLOP count overcomes the kernel overhead and non-attention fixed costs, resulting in the $1.09\times$--$1.27\times$ end-to-end prefill speedups. At longer and longer context, quadratic dense attention dominates prefill compute and the theoretical comparison ratio continues to upper-bound the empirical prefill speedup.

\subsection{Approximation Error and Expressivity Bounds}

Because Bigger Bird~v3 is applied training-free to pretrained dense checkpoints, its primary theoretical objective is approximating dense causal attention rather than learning new multi-hop routing weights from scratch. For any query $i$ in a given layer and head, let $A_i = (A_{i,1}, \dots, A_{i,i})$ denote the exact dense causal softmax weights, $S_i \subseteq \{1, \dots, i\}$ the retained key set, and $\epsilon_i = \sum_{j \notin S_i} A_{i,j} \in [0, 1)$ the omitted attention probability mass. Renormalizing softmax over $S_i$ gives sparse weights $\tilde{A}_{i,j} = A_{i,j} / (1 - \epsilon_i)$ for $j \in S_i$ and $0$ otherwise, whose $\ell_1$ deviation from dense attention is exact: 

\begin{equation}
\lVert A_i - \tilde{A}_i \rVert_1
  = \sum_{j \in S_i} \left(\frac{A_{i,j}}{1 - \epsilon_i} - A_{i,j}\right)
  + \sum_{j \notin S_i} A_{i,j}
  = \frac{\epsilon_i}{1 - \epsilon_i}(1 - \epsilon_i) + \epsilon_i
  = 2\epsilon_i.
\label{eq:l1-mass-error}
\end{equation}

Consequently, the Euclidean error between the dense output $o_i = \sum_{j=1}^i A_{i,j} v_j$ and the sparse output $\tilde{o}_i = \sum_{j \in S_i} \tilde{A}_{i,j} v_j$ satisfies

\begin{equation}
\lVert o_i - \tilde{o}_i \rVert_2
  = \left\lVert \sum_{j=1}^i (A_{i,j} - \tilde{A}_{i,j}) v_j \right\rVert_2
  \le 2\epsilon_i \max_{1 \le j \le i} \lVert v_j \rVert_2.
\label{eq:l2-output-error}
\end{equation}

This bound motivates each architectural component of Bigger Bird~v3 directly: the $64$ attention sinks and $16{,}384$-token local window capture the structurally persistent attention mass associated with initial sink tokens~\cite{xiao2024streamingllm} and local recency; the per-head content router targets up to $512$ high-mass distant columns to suppress the remaining tail mass $\epsilon_i$; and the exact $256$-token tail enforces $\epsilon_i = 0$ identically at the end of the prompt, eliminating single-layer attention truncation error where answer aggregation occurs.

Original BigBird~\cite{zaheer2020bigbird} established three foundational theoretical results for sparse Transformers: universal approximation of continuous sequence-to-sequence functions via a bidirectional star graph, Turing completeness under arbitrary precision, and a depth lower bound under the Orthogonal Vectors Conjecture (OVC). We apply and discuss these frameworks in relation to BiggerBird~v3 below:

\textbf{Universal approximation and causal reachability:}
The Universal Approximation Theorem established in the original BigBird~\cite{zaheer2020bigbird} applies to bidirectional encoders, while Bigger Bird~v3 operates under a causal decoder mask where early sink tokens cannot attend to future positions. Restricting the target function class to continuous \emph{causal} sequence-to-sequence mappings, v3's $16{,}384$-token sliding window still guarantees full causal reachability between any $j \le i$ in at most $\lceil (i - j) / W \rceil$ hops, satisfying the graph-connectivity requirement for idealized causal universality under unbounded depth, trainable weights, and arbitrary precision \cite{yun2020connections}. In practice, however, universal approximation proofs are existential and rely on specialized multi-hop relay weights. Because v3 runs training-free on frozen, fixed-depth weights optimized for single-hop dense attention, multi-hop propagation through the window alone is unreliable, a theoretical gap confirmed empirically by Table~\ref{tab:v3-ablate} where removing the exact tail (\texttt{exact\_tail = 0}) drops $\sim 127\text{K}$ needle retrieval from $6/6$ to $2/6$, while restoring \texttt{exact\_tail = 256} recovers $6/6$.

\textbf{Depth lower bound (OVC):}
Outside the final $256$ tokens, Bigger Bird~v3 computes $\mathcal{O}(n)$ inner products per layer. The original BigBird's lower bound~\cite{zaheer2020bigbird} therefore continues to hold for the prefill body: under the Orthogonal Vectors Conjecture, solving simultaneous all-pairs global comparison tasks (such as finding the furthest vector for every token in parallel) requires $\widetilde{\Omega}(n^{1-o(1)})$ layers. For inference queries positioned at the end of the prompt or during generation, however, v3's exact $256$-token tail and exact decode over the retained KV cache bypass this all-pairs bottleneck in $\mathcal{O}(1)$ layers.

\textbf{Turing completeness, precision, and KV eviction:}
Unlike BigBird's sparse-decoder Turing completeness construction, which required an explicit causal sparse graph and $\mathcal{O}(\sqrt{i})$ intermediate scratchpad steps per tape lookup~\cite{zaheer2020bigbird}, Bigger Bird~v3 retains exact dense attention during autoregressive decode. Under the standard idealized assumptions of unbounded memory and arbitrary numerical precision, exact decode over the full KV cache trivially preserves Turing completeness \cite{perez2019turing}. In deployment, however, two practical constraints strictly reduce the system to a bounded finite-state automaton: BF16 precision bounds addressable state, and enabling optional KV eviction (\texttt{bb\_evict64k} in Table~\ref{tab:r1-evict} or fixed $16\text{K}$--$128\text{K}$ budgets in Experiment~21, Table~\ref{tab:1m-exp21}) discards prefix entries and removes access to an unbounded tape history.


\section{Limitations}
\label{sec:limitations}

The primary evaluation is limited to R1-Distill-Llama-8B on a single H100 GPU, with quality measured through first-token multiple-choice accuracy rather than full autoregressive generation or downstream task completion. Moreover, dense accuracy on Code.Debug and LongBench v2 was close to chance, limiting conclusions on reasoning-quality preservation. Furthermore, nonsignificant accuracy differences do not establish equivalence, and prompts exceeding 130k tokens were middle-truncated, potentially removing task-relevant information. 

The R1 timing protocol used short warmups and a single timed forward per example, limiting the reliability of latency estimates. The exact upstream R1 weight revision was also not recorded in Modal metadata, reducing reproducibility. The near 1M token LongCodeQA pilot covers only 18 of 50 prompts under a separate no-KV-cache protocol, preventing conclusions on end-to-end generation performance at this scale. Similarly, the KV-eviction experiments do not establish long-generation quality, as eviction occurs after first-token logits are computed, nor do they demonstrate peak-memory savings under the evaluated protocol. 

Selecting the $512$ distant routed positions from the mean causal attention probabilities of the final $64$ prompt queries relies on a late-query proxy assumption: that the keys most salient to the end of the prompt cover the long-range context needed by earlier prefill tokens. This assumption aligns well with prompts whose question or instruction appears at the end (the same structural property assumed by our middle-truncation protocol), but can weaken on instruction-first prompts where the final $64$ tokens consist solely of document text. Moreover, because our primary R1 evaluation scores first-token multiple-choice logits without chain-of-thought generation, it does not measure routing drift during extended multi-step decoding, where newly generated reasoning steps may depend on distant facts whose prefix KV representations were constructed without attending to one another during prefill.

Broadcasting routing indices chosen by the final $64$ prompt queries to earlier prefill tokens is causally valid during inference prefill, where the entire prompt is fixed prior to generation, but cannot be applied unmodified to parallel autoregressive training. Allowing an intermediate token at position $i < n - 64$ to select keys using attention scores from the final $64$ tokens would leak future information backward into next-token predictions at position $i$. Extending Bigger Bird to the sparse-attention training regime noted in Section~\ref{sec:conclusions} will therefore require strictly causal per-query or per-block routing mechanisms.

Finally, the individual contributions of Bigger Bird v3's attention components have not been fully isolated through ablations. Published sparse-attention methods (DSA, LSA, and ZHA) have not been evaluated under the matched protocol, leaving comparative performance unresolved. The earlier IMDb and kernel experiments also use unequal training budgets, incomplete metrics, and single-run estimates, and should therefore be interpreted as developmental evidence rather than definitive architectural comparisons. 

\section{Conclusions and Recommendations}
\label{sec:conclusions}
Bigger Bird v3 is a context-aware sparse-prefill method for long-context decoder inference. On the matched H100 R1 suite, it remains close to dense multiple-choice accuracy while improving prefill on longer prompts. Next steps are completing the remaining 1M LongCodeQA prompts, longer-generation and
eviction quality tests, sparse-attention training, and matched comparisons against recent inference-time sparse methods.

Together, these results suggest that attention routing can substantially reduce long-context prefill cost while keeping pooled accuracy close to dense attention, though the accuracy gap and limited 1M-prompt pilot leave quality at larger contexts to be tested.


\clearpage
\addcontentsline{toc}{section}{References}
\bibliographystyle{IEEEtranN}
\bibliography{references}

\clearpage
\appendix
\input{appendix}

\end{document}